\begin{table}[htbp]
\centering
\footnotesize
\caption{Kuramoto K-sweep results by Pi initial frequency condition.}
\label{tab:k_sweep_condition}
\begin{tabular}{@{}lrrrrrrr@{}}
\toprule
\textbf{Pi} & \textbf{K} & \textbf{Trials} & \textbf{FCR} & \textbf{5 s MAE} & \textbf{Virt. std} & \textbf{Pi std} & \textbf{Score} \\
\midrule
1.2 Hz & 2.5 & 2 & 0.877 & 0.687 & 0.263 & 0.251 & 1.029 \\
1.5 Hz & 2.5 & 2 & 0.935 & 0.315 & 0.196 & 0.258 & 0.632 \\
2.5 Hz & 2.5 & 2 & 0.721 & 0.293 & 0.180 & 0.261 & 0.881 \\
1.2 Hz & 3.0 & 2 & 0.952 & 0.560 & 0.355 & 0.324 & 2.140 \\
1.5 Hz & 3.0 & 2 & 0.900 & 0.213 & 0.155 & 0.272 & 1.103 \\
2.5 Hz & 3.0 & 2 & 0.740 & 0.186 & 0.086 & 0.234 & 1.946 \\
1.2 Hz & 3.5 & 2 & 0.932 & 0.618 & 0.464 & 0.360 & 2.819 \\
1.5 Hz & 3.5 & 2 & 0.918 & 0.151 & 0.093 & 0.158 & 0.977 \\
2.5 Hz & 3.5 & 2 & 0.767 & 0.318 & 0.229 & 0.340 & 3.698 \\
1.2 Hz & 4.0 & 2 & 0.934 & 0.476 & 0.350 & 0.362 & 2.433 \\
1.5 Hz & 4.0 & 2 & 0.844 & 0.370 & 0.291 & 0.371 & 2.104 \\
2.5 Hz & 4.0 & 2 & 0.757 & 0.352 & 0.178 & 0.370 & 3.572 \\
1.2 Hz & 4.5 & 2 & 0.931 & 0.315 & 0.063 & 0.389 & 1.911 \\
1.5 Hz & 4.5 & 2 & 0.938 & 0.376 & 0.326 & 0.260 & 2.159 \\
2.5 Hz & 4.5 & 2 & 0.736 & 0.306 & 0.150 & 0.383 & 3.311 \\
\bottomrule
\end{tabular}
\end{table}
